\section{API Design and Endpoints}
\label{sec:api_design}

The communication between the frontend client and the backend server is
established through a structured RESTful API. This design ensures that the
system remains stateless and scalable, utilizing standard HTTP methods to
perform operations on resources. The Node.js application server acts as the
centralized gateway, exposing a unified set of endpoints that handle both
standard clinical data and custom graph logic.

\subsection{RESTful Structure and Formats}
All API interactions strictly adhere to the Javascript Object Notation (JSON)
format for both request bodies and response payloads. The system utilizes
standard HTTP verbs to define the nature of the operation. GET requests are
used exclusively for retrieving data without side effects, while POST requests
handle the creation of new resources. PUT methods are employed for updating
existing records, and DELETE methods remove specific resources when permitted.

Responses from the server are standardized to ensure predictability for the
client application. Successful operations return a 200 OK status code
accompanied by the requested data object. In cases of resource creation, a 201
Created status is returned. Error handling is centralized, ensuring that any
failure creates a consistent JSON error object containing a human-readable
message and a specific error code, rather than exposing raw server stack
traces.

\subsection{FHIR-Compliant Endpoints}
A significant portion of the API is dedicated to proxying standard FHIR
resources. Endpoints such as \texttt{/api/patients} and
\texttt{/api/encounters} act as a validation and caching layer sitting in front
of the HAPI FHIR server. When the frontend requests patient details, the
Node.js server intercepts this call, validates the parameters, and then
forwards the request to the underlying Java server. The response is strictly
strictly compliant with the FHIR R5 standard, ensuring that the data structure
remains interoperable with other external healthcare systems.

\subsection{Custom History Graph Endpoints}
Beyond standard FHIR resources, the API exposes specialized endpoints designed
to support the Patient History Graph visualization. The primary endpoint,
\texttt{/api/historyGraph}, does not return a simple list of records. Instead,
it aggregates data from multiple FHIR resources and transforms them into a
graph structure containing nodes and relational edges. This allows the frontend
to render the complex temporal relationships between medical events without
needing to perform heavy data processing in the browser.

\begin{table}[H]
    \centering
    \begin{tabular}{|l|l|p{5cm}|}
        \hline
        \textbf{Method} & \textbf{Endpoint}              & \textbf{Description}                                           \\
        \hline
        \hline
        GET             & \texttt{/api/patients}         & Retrieve list of patients with pagination support              \\
        \hline
        POST            & \texttt{/api/patients}         & Register a new patient with FHIR validation                    \\
        \hline
        GET             & \texttt{/api/encounters/:id}   & Get specific encounter details including clinical observations \\
        \hline
        GET             & \texttt{/api/historyGraph/:id} & Get patient history graph nodes and causal relationships       \\
        \hline
        POST            & \texttt{/auth/login}           & Authenticate user and obtain JWT token via Keycloak            \\
        \hline
    \end{tabular}
    \caption{Key REST API Endpoints provided by the Node.js Backend.}
    \label{tab:api_endpoints}
\end{table}